% ==============================================================================
% PLANTILLA DE PRESENTACIÓN INSTITUCIONAL - UANL
% Universidad Autónoma de Nuevo León | Facultad de Agronomía / FIME
% ==============================================================================
% Inspirada en las directrices de identidad de AulaTeX (estilo UnADM / UANL),
% adaptada a la paleta institucional (Azul y Oro UANL, Verde Agronomía)
% y al flujo de defensa síncrona de reportes y Productos Integradores (PIA).
%
% Uso:
% 1. Duplica o importa este archivo para presentaciones de la UANL.
% 2. Actualiza los metadatos editables: título, actividad, asignatura, docente,
%    matrícula, fecha y bibliografía.
% 3. Compilar con:
%    powershell -File .\aulaTEX\latexmk-build.ps1 .\UANL\presentacion-uanl.tex
% ==============================================================================

\documentclass[
  spanish,
  aspectratio=169,
  xcolor={dvipsnames,table}
]{beamer}

% Proporción panorámica 16:9 (19.2 cm x 10.8 cm)
\geometry{paperwidth=19.2cm,paperheight=10.8cm}

\usepackage{ifxetex}
\ifxetex
  \usepackage{fontspec}
  \setmainfont{Arial}
  \setsansfont{Arial}
\else
  \usepackage[utf8]{inputenc}
  \usepackage[T1]{fontenc}
  \usepackage[scaled=.96]{helvet}
  \renewcommand{\familydefault}{\sfdefault}
\fi

\usepackage[spanish,es-tabla,es-noshorthands]{babel}
\usepackage{array}
\usepackage{booktabs}
\usepackage{graphicx}
\usepackage{ragged2e}
\usepackage{tabularx}
\usepackage{amsmath,amssymb}
\usepackage{tikz}
\usepackage{hyperref}
\usetikzlibrary{arrows.meta,calc,fit,matrix,positioning,shapes.geometric,shadows.blur}

\newcolumntype{Y}{>{\RaggedRight\arraybackslash}X}

% ------------------------------------------------------------------------------
% DATOS Y METADATOS EDITABLES
% ------------------------------------------------------------------------------
\providecommand{\studentname}{Carlos Mauricio Carvajal Coronado}
\providecommand{\studentshort}{C. M. Carvajal Coronado}
\providecommand{\studentid}{1977875}
\providecommand{\teamname}{Célula Base / Individual}
\providecommand{\universityname}{Universidad Autónoma de Nuevo León}
\providecommand{\facultyname}{Facultad de Agronomía}
\providecommand{\departmentname}{Ingeniero Agrónomo}
\providecommand{\coursename}{Biología Celular}
\providecommand{\coursecode}{REF\#261013755061}
\providecommand{\activitytitle}{Plantilla de Presentación Institucional UANL}
\providecommand{\activitysubtitle}{Defensa de Evidencias de Aprendizaje y PIA}
\providecommand{\teachingfigure}{Dr. Miguel Ángel Dector Carrillo}
\providecommand{\deliverydate}{\today}
\providecommand{\locationname}{General Escobedo, N.L.}
\providecommand{\departmentlogo}{img/departamentos/logo-agronomia}

% ------------------------------------------------------------------------------
% PALETA CROMÁTICA INSTITUCIONAL UANL Y FACULTAD DE AGRONOMÍA
% ------------------------------------------------------------------------------
\definecolor{uanlBlueDark}{HTML}{002855}   % Azul UANL profundo
\definecolor{uanlBlue}{HTML}{00376C}       % Azul institucional UANL
\definecolor{uanlGold}{HTML}{D4AF37}       % Oro / Dorado universitario UANL
\definecolor{uanlGoldLight}{HTML}{F0A500}  % Oro de acento
\definecolor{uanlGreenAgro}{HTML}{1B5E20}  % Verde Facultad de Agronomía
\definecolor{uanlGreenLight}{HTML}{2E7D32} % Verde follaje
\definecolor{uanlPaper}{HTML}{F6F8F7}      % Fondo tenue para cajas
\definecolor{uanlInk}{HTML}{1A2421}        % Tipografía principal
\definecolor{uanlMuted}{HTML}{546E66}      % Texto atenuado

\mode<presentation>{
  \usetheme{default}
  \usefonttheme{professionalfonts}
  \setbeamertemplate{navigation symbols}{}
  \setbeamertemplate{blocks}[rounded][shadow=false]
  \setbeamercovered{invisible}
}
\setbeamersize{text margin left=0.65cm,text margin right=0.65cm}

\hypersetup{
  colorlinks=true,
  urlcolor=uanlBlue,
  linkcolor=uanlBlue
}

\setbeamercolor{background canvas}{bg=white}
\setbeamercolor{normal text}{fg=uanlInk,bg=white}
\setbeamercolor{structure}{fg=uanlBlue}
\setbeamercolor{frametitle}{fg=white,bg=uanlBlue}
\setbeamercolor{title}{fg=white,bg=uanlBlue}
\setbeamercolor{subtitle}{fg=white}
\setbeamercolor{block title}{fg=white,bg=uanlBlue}
\setbeamercolor{block body}{fg=uanlInk,bg=uanlPaper}
\setbeamercolor{block title alerted}{fg=white,bg=uanlGold}
\setbeamercolor{block body alerted}{fg=uanlInk,bg=uanlPaper}
\setbeamercolor{block title example}{fg=white,bg=uanlGreenAgro}
\setbeamercolor{block body example}{fg=uanlInk,bg=uanlPaper}
\setbeamercolor{alerted text}{fg=uanlGoldLight}
\setbeamerfont{title}{size=\LARGE,series=\bfseries}
\setbeamerfont{frametitle}{size=\large,series=\bfseries}
\setbeamerfont{block title}{series=\bfseries}
\setbeamertemplate{itemize item}{\textcolor{uanlGreenAgro}{\large$\blacktriangleright$}}
\setbeamertemplate{itemize subitem}{\textcolor{uanlGold}{\scriptsize$\blacksquare$}}
\setbeamertemplate{enumerate item}{\textcolor{uanlBlue}{\insertenumlabel.}}

% ------------------------------------------------------------------------------
% ENCABEZADO INSTITUCIONAL
% ------------------------------------------------------------------------------
\setbeamertemplate{frametitle}{
  \nointerlineskip
  \begin{beamercolorbox}[wd=\textwidth,ht=1.05cm,dp=0.22cm,leftskip=0.40cm,rightskip=0.25cm]{frametitle}
    \begin{minipage}[c]{0.74\textwidth}
      \insertframetitle\par
      \ifx\insertframesubtitle\@empty\else%
        {\fontsize{9}{11}\selectfont\color{uanlGoldLight}\insertframesubtitle}%
      \fi
    \end{minipage}
    \hfill
    \raisebox{-0.1cm}{\includegraphics[height=0.65cm]{\departmentlogo}}
  \end{beamercolorbox}
  \vspace{-0.04cm}
  {\color{uanlGold}\rule{\textwidth}{1.4pt}}
}

% ------------------------------------------------------------------------------
% PIE DE PÁGINA INSTITUCIONAL
% ------------------------------------------------------------------------------
\setbeamertemplate{footline}{
  \leavevmode%
  \hbox{%
    \begin{beamercolorbox}[wd=.33\paperwidth,ht=2.8ex,dp=1.1ex,leftskip=1.2em]{author in head/foot}%
      \color{white}\studentshort\ \textbar\ \studentid
    \end{beamercolorbox}%
    \begin{beamercolorbox}[wd=.42\paperwidth,ht=2.8ex,dp=1.1ex,center]{title in head/foot}%
      \color{white}\textbf{\coursename}\ (\facultyname)
    \end{beamercolorbox}%
    \begin{beamercolorbox}[wd=.25\paperwidth,ht=2.8ex,dp=1.1ex,rightskip=1.2em plus 1fil]{date in head/foot}%
      \color{white}\coursecode\hfill\insertframenumber/\inserttotalframenumber
    \end{beamercolorbox}%
  }%
  \vskip0pt%
}
\setbeamercolor{author in head/foot}{bg=uanlBlueDark}
\setbeamercolor{title in head/foot}{bg=uanlGreenAgro}
\setbeamercolor{date in head/foot}{bg=uanlBlueDark}

% ------------------------------------------------------------------------------
% BARRA DE ACENTO INFERIOR
% ------------------------------------------------------------------------------
\newcommand{\accentbar}{%
  \begin{tikzpicture}[remember picture,overlay]
    \path[use as bounding box] (0,0) rectangle (0,0);
    \fill[uanlGreenAgro] (current page.south west) rectangle ([xshift=0.55\paperwidth,yshift=0.08cm]current page.south west);
    \fill[uanlGold] ([xshift=0.55\paperwidth]current page.south west) rectangle ([xshift=\paperwidth,yshift=0.08cm]current page.south west);
  \end{tikzpicture}
}

% ------------------------------------------------------------------------------
% SEPARADOR DE SECCIONES DE ALTO IMPACTO VISUAL
% ------------------------------------------------------------------------------
\newcommand{\uanlsectionframe}[1]{%
  \begin{frame}[plain]
    \begin{tikzpicture}[remember picture,overlay]
      \path[use as bounding box] (0,0) rectangle (0,0);
      \fill[uanlBlueDark] (current page.south west) rectangle (current page.north east);
      \node[opacity=0.08,anchor=east] at ([xshift=-0.5cm]current page.east) {\includegraphics[width=0.45\paperwidth]{\departmentlogo}};
      \fill[uanlGreenAgro] ([xshift=1.05cm,yshift=2.25cm]current page.south west) rectangle ([xshift=1.25cm,yshift=6.50cm]current page.south west);
      \fill[uanlGold] ([xshift=1.45cm,yshift=2.25cm]current page.south west) rectangle ([xshift=1.65cm,yshift=5.35cm]current page.south west);
      \node[
        anchor=west,
        align=left,
        text=white,
        text width=0.62\paperwidth,
        font=\fontsize{24}{29}\selectfont\bfseries
      ] at ([xshift=2.10cm,yshift=0.60cm]current page.west) {#1};
      \node[
        anchor=west,
        align=left,
        text=uanlGoldLight,
        font=\large\bfseries
      ] at ([xshift=2.10cm,yshift=-0.60cm]current page.west) {\facultyname\ \textbar\ \universityname};
    \end{tikzpicture}
  \end{frame}
}

\newcommand{\uanlsection}[1]{%
  \section{#1}
  \uanlsectionframe{#1}
}

% ------------------------------------------------------------------------------
% MACRO PARA PORTADA INSTITUCIONAL
% ------------------------------------------------------------------------------
\newcommand{\uanlportada}{%
  \begin{frame}[plain]
    \begin{tikzpicture}[remember picture,overlay]
      \path[use as bounding box] (0,0) rectangle (0,0);
      \fill[uanlBlueDark,opacity=0.98] (current page.south west) rectangle ([xshift=0.62\paperwidth]current page.north west);
      \fill[uanlGreenAgro,opacity=0.94] ([xshift=0.62\paperwidth]current page.south west) rectangle (current page.north east);
      \node[anchor=center,opacity=0.08] at (current page.center) {\includegraphics[width=0.75\paperwidth]{\departmentlogo}};
      \node[anchor=north west] at ([xshift=0.65cm,yshift=-0.55cm]current page.north west) {\includegraphics[width=3.2cm]{\departmentlogo}};
      \fill[uanlGold] ([xshift=0.64\paperwidth,yshift=1.25cm]current page.south west) rectangle ([xshift=0.95\paperwidth,yshift=1.38cm]current page.south west);
      \fill[white,opacity=0.18] ([xshift=0.64\paperwidth,yshift=0.98cm]current page.south west) rectangle ([xshift=0.86\paperwidth,yshift=1.08cm]current page.south west);
    \end{tikzpicture}

    \vspace*{1.75cm}
    \begin{minipage}{0.58\paperwidth}
      {\color{white}\fontsize{21}{25}\selectfont\bfseries\raggedright \activitytitle\par}
      \vspace{0.20cm}
      {\color{uanlGoldLight}\large\bfseries \activitysubtitle}\\[0.28cm]
      {\color{uanlGold}\rule{0.85\linewidth}{1.4pt}}\\[0.35cm]
      {\color{white}\textbf{Estudiante:} \studentname}\\
      {\color{white!85}\small\textbf{Matrícula:} \studentid \quad\textperiodcentered\quad \textbf{Equipo:} \teamname}
    \end{minipage}
    \hfill
    \begin{minipage}{0.33\paperwidth}
      \raggedleft
      {\color{white}\Large\bfseries UANL}\\[0.12cm]
      {\color{white!90}\footnotesize \universityname}\\[0.35cm]
      {\color{white}\small\bfseries \facultyname}\\
      {\color{white!85}\small \departmentname}\\[0.35cm]
      {\color{uanlGoldLight}\small\textbf{\coursename}}\\
      {\color{white!75}\footnotesize Clave: \coursecode}\\[0.35cm]
      {\color{white!80}\footnotesize \textbf{Docente:}\\ \teachingfigure}\\[0.25cm]
      {\color{white!70}\scriptsize \locationname\ \textperiodcentered\ \deliverydate}
    \end{minipage}
  \end{frame}
}

% ==============================================================================
% INICIO DE LA PRESENTACIÓN
% ==============================================================================
\begin{document}

% 1. PORTADA
\uanlportada

% 2. OBJETIVO Y ENFOQUE INSTITUCIONAL
\begin{frame}{Objetivo, competencias y alineación institucional}
  \begin{columns}[T]
    \begin{column}{0.48\textwidth}
      \begin{block}{Objetivo formativo}
        Estructurar una defensa oral concisa, rigurosa y visualmente orientada que exponga los fundamentos biofísicos, metodológicos y analíticos del trabajo académico desarrollado en la UANL.
      \end{block}
      \begin{alertblock}{Criterios de acreditación UANL}
        \begin{itemize}
          \item Rigor conceptual sin ambigüedades técnicas.
          \item Trazabilidad de fuentes formales en APA 7.
          \item Capacidad de síntesis en tiempo asignado.
        \end{itemize}
      \end{alertblock}
    \end{column}
    \begin{column}{0.48\textwidth}
      \begin{exampleblock}{Competencias del Modelo Educativo}
        \begin{itemize}
          \item \textbf{Pensamiento crítico:} Explicación de relaciones causa-efecto en fenómenos biológicos o de ingeniería.
          \item \textbf{Comunicación científica:} Dominio de terminología especializada y esquemas vectoriales.
          \item \textbf{Transferencia:} Vinculación de principios teóricos con la resolución de problemas agropecuarios o industriales.
        \end{itemize}
      \end{exampleblock}
    \end{column}
  \end{columns}
  \accentbar
\end{frame}

% 3. SECCIÓN 1: FUNDAMENTOS TEÓRICOS
\uanlsection{1. Fundamentos Científicos y Bases Metodológicas}

\begin{frame}{Ejes temáticos y dimensiones analíticas}
  \begin{columns}[T]
    \begin{column}{0.48\textwidth}
      \begin{block}{Dimensión estructural y bioenergética}
        \begin{itemize}
          \item Análisis de la ultraestructura celular y organización subcelular.
          \item Gradientes electroquímicos y transporte transmembranal sostenido por bombas de protones ($H^+$-ATPasa).
          \item Homeostasis metabólica y modulación frente a estrés ambiental.
        \end{itemize}
      \end{block}
    \end{column}
    \begin{column}{0.48\textwidth}
      \begin{exampleblock}{Metodología y sistematización}
        \begin{itemize}
          \item Selección y depuración de literatura científica en bases indexadas (Scopus, WoS, PubMed).
          \item Análisis comparativo transversal entre modelos vegetales y animales.
          \item Integración de evidencia cuantitativa y experimental verificada.
        \end{itemize}
      \end{exampleblock}
    \end{column}
  \end{columns}
  \accentbar
\end{frame}

% 4. SECCIÓN 2: EVIDENCIA VISUAL
\uanlsection{2. Evidencia Visual y Modelado de Procesos}

\begin{frame}{Arquitectura de flujo analítico y toma de decisiones}
  \centering
  \begin{tikzpicture}[
    node distance=1.4cm,
    proc/.style={
      rectangle,
      draw=uanlBlue,
      fill=uanlPaper,
      thick,
      rounded corners=4pt,
      text width=3.8cm,
      align=center,
      font=\small\sffamily,
      minimum height=1.3cm
    },
    agro/.style={
      rectangle,
      draw=uanlGreenAgro,
      fill=white,
      thick,
      rounded corners=4pt,
      text width=3.8cm,
      align=center,
      font=\small\sffamily\bfseries,
      minimum height=1.3cm
    },
    conn/.style={
      -Stealth,
      thick,
      draw=uanlBlueDark
    }
  ]
    \node[proc] (p1) {\textbf{Señal Ambiental}\\\footnotesize Detección de estímulos (luz, sequía, temperatura)};
    \node[proc, right=1.2cm of p1] (p2) {\textbf{Cascada Intracelular}\\\footnotesize Segundos mensajeros y fosforilación ($Ca^{2+}$, SnRK2.6)};
    \node[agro, right=1.2cm of p2] (p3) {\textbf{Respuesta Fisiológica}\\\footnotesize Modulación estomática y eficiencia en uso del agua (WUE)};

    \path[conn] (p1) -- (p2);
    \path[conn] (p2) -- (p3);
  \end{tikzpicture}
  \vspace{0.8cm}
  \begin{block}{Interpretación del modelo}
    La integración biofísica traduce gradientes osmóticos en respuestas morfológicas macroscópicas, permitiendo evaluar la tolerancia vegetal frente al estrés hídrico en sistemas de producción agrícola.
  \end{block}
  \accentbar
\end{frame}

% 5. CONCLUSIONES Y TRANSFERENCIA
\uanlsection{3. Conclusiones y Transferencia Profesional}

\begin{frame}{Síntesis y aplicación práctica en agronomía e industria}
  \begin{columns}[T]
    \begin{column}{0.48\textwidth}
      \begin{block}{Principios consolidados}
        \begin{enumerate}
          \item La estructura celular determina la competencia fisiológica y la resiliencia productiva del cultivo.
          \item El aislamiento simplástico previene la disipación pasiva de potenciales osmóticos foliares.
          \item La señalización hormonal rige la respuesta adaptativa ante fluctuaciones climáticas.
        \end{enumerate}
      \end{block}
    \end{column}
    \begin{column}{0.48\textwidth}
      \begin{exampleblock}{Transferencia agronómica directa}
        \begin{itemize}
          \item \textbf{Manejo de riego de precisión:} Programación hídrica basada en conductancia estomática.
          \item \textbf{Mejoramiento genético:} Selección de variedades con alta eficiencia en el uso del agua (WUE).
          \item \textbf{Biotecnología:} Optimización de bioestimulantes osmoprotectores foliares.
        \end{itemize}
      \end{exampleblock}
    \end{column}
  \end{columns}
  \accentbar
\end{frame}

% 6. REFERENCIAS BIBLIOGRÁFICAS
\begin{frame}{Referencias bibliográficas formales (APA 7ma edición)}
  \begin{columns}[T]
    \begin{column}{0.48\textwidth}
      \footnotesize
      \textbf{Normativa y directrices institucionales UANL:}
      \begin{itemize}
        \item Universidad Autónoma de Nuevo León. (2022). \emph{Modelo Educativo de la UANL}. Monterrey: Casa Universitaria del Libro.
        \item Universidad Autónoma de Nuevo León. (2019). \emph{Plan de Desarrollo Institucional: Visión 2030 UANL}. San Nicolás de los Garza: Rectoría UANL.
        \item H. Consejo Universitario UANL. (2021). \emph{Estatuto General de la Universidad Autónoma de Nuevo León}.
      \end{itemize}
    \end{column}
    \begin{column}{0.48\textwidth}
      \footnotesize
      \textbf{Tratados y literatura científica especializada:}
      \begin{itemize}
        \item Alberts, B., et al. (2022). \emph{Molecular Biology of the Cell} (7th ed.). W. W. Norton \& Company.
        \item Nelson, D. L., \& Cox, M. M. (2021). \emph{Lehninger Principles of Biochemistry} (8th ed.). W. H. Freeman.
        \item Taiz, L., et al. (2015). \emph{Plant Physiology and Development} (6th ed.). Sinauer Associates / Oxford University Press.
        \item Hernández-Sampieri, R., \& Mendoza, C. P. (2018). \emph{Metodología de la investigación}. McGraw-Hill.
      \end{itemize}
    \end{column}
  \end{columns}
  \accentbar
\end{frame}

\end{document}
